\section{Introducción al Dossier de Diseño Mecánico e Industrial}

Este dossier presenta la especificación mecánica e industrial del collar de geocercamiento ganadero DR-VF-01. Se detallan las estrategias de estanqueidad bajo norma IP67/IP68, la asignación de tolerancias dimensionales y geométricas, la integración de las tres antenas (LoRa, Wi-Fi/BLE y GNSS) y la disposición estructural del panel solar y la carcasa.

\section{Especificaciones de Sellado y Estanqueidad IP67}

\subsection{Requisitos de Protección Ambiental}

El gabinete del collar DR-VF-01 está diseñado para cumplir con el grado de protección IP67 (según norma IEC 60529), garantizando resistencia ante las condiciones severas del entorno de pastoreo:

\begin{itemize}[leftmargin=*]
    \item \textbf{Protección contra Polvo y Partículas (Grado 6)}: Sellado hermético total contra la penetración de material particulado fino.
    \item \textbf{Protección contra Inmersión en Agua (Grado 7)}: Inmunidad ante inmersión temporal a $1.0\,\text{m}$ de profundidad durante $30\,\text{min}$.
\end{itemize}

\subsection{Diseño Estructural de Sellado}

\subsubsection{Carcasa Principal y Juntas Óptimas}

La carcasa está inyectada en policarbonato reforzado con $20\%$ de fibra de vidrio (\text{PC+GF20\%}), asegurando un espesor de pared estructural de $2.5\,\text{mm}$. La estanqueidad entre tapas se logra mediante un canal perimetral con junta tórica (\textit{O-ring}) de elastómero EPDM, calculada para una compresión nominal del $20\text{--}25\%$.

\subsubsection{Protección de Interfaz de Carga USB-C}

El puerto USB-C sirve para la programación por USB nativo y para la carga mediante el cargador TP4056, y debe sellarse para cumplir IP67. El tipo de sello (tapón ciego, tapa u otro) y el MPN del conector USB-C están \textbf{POR DECIDIR}; hasta entonces este dossier no fija la solución de sellado de este puerto.

\subsection{Procedimientos de Verificación de Estanqueidad}

\subsubsection{Ensayos de Cámara de Polvo}

\begin{itemize}[leftmargin=*]
    \item \textbf{Polvo de Ensayo}: Talco fino con granulometría $\le 75\,\mu\text{m}$.
    \item \textbf{Duración}: $8\,\text{horas}$ en depresión de $2.0\,\text{kPa}$.
    \item \textbf{Criterio de Aprobación}: Ausencia total de partículas sobre la tarjeta madre (PCBA).
\end{itemize}

\subsubsection{Ensayos de Inmersión Hidrostática}

\begin{itemize}[leftmargin=*]
    \item \textbf{Profundidad de Inmersión}: $1.0\,\text{m}$ desde la superficie del líquido.
    \item \textbf{Duración del Ensayo}: $30\,\text{minutos}$ continuos.
    \item \textbf{Temperatura de Agua}: $20\,^\circ\text{C} \pm 5\,^\circ\text{C}$.
    \item \textbf{Criterio de Aprobación}: Incremento de masa $\Delta m = 0.00\,\text{g}$ y verificación dieléctrica posterior.
\end{itemize}

\section{Tolerancias Dimensionales y Ajustes}

\subsection{Clasificación de Tolerancias}

\begin{description}[leftmargin=*, style=nextline, font=\bfseries\ttfamily]
    \item[Clase A (Crítica)] $\pm 0.05\,\text{mm}$ aplicable a alojamientos de junta \textit{O-ring}, guías de pines y planos de sellado.
    \item[Clase B (Importante)] $\pm 0.10\,\text{mm}$ para apoyos del PCBA, torretas de fijación y encajes de batería.
    \item[Clase C (General)] $\pm 0.20\,\text{mm}$ para dimensiones exteriores globales de la carcasa.
\end{description}

\subsection{Tolerancias Críticas de la Carcasa}

\begin{table}[H]
\centering
\caption{Tolerancias dimensionales de la carcasa principal del collar DR-VF-01}
\label{tab:tolerancias_carcasa}
\vspace{2mm}
\begin{tabularx}{\textwidth}{X r r c}
\toprule
\rowcolor{categorybg}
\textbf{Dimensión / Cota} & \textbf{Valor Nominal} & \textbf{Tolerancia} & \textbf{Clase ISO} \\
\midrule
Longitud Total & $120.0\,\text{mm}$ & $\pm 0.20\,\text{mm}$ & Clase C \\
\rowcolor{rowalt}
Anchura Total & $80.0\,\text{mm}$ & $\pm 0.20\,\text{mm}$ & Clase C \\
Altura Total & $25.0\,\text{mm}$ & $\pm 0.20\,\text{mm}$ & Clase C \\
\rowcolor{rowalt}
Espesor de Pared Estructural & $2.5\,\text{mm}$ & $\pm 0.10\,\text{mm}$ & Clase B \\
Ancho de Canal de Sellado & $1.8\,\text{mm}$ & $\pm 0.05\,\text{mm}$ & Clase A \\
\bottomrule
\end{tabularx}
\end{table}

\subsection{Tolerancias Geométricas y de Acabado de Superficie}

\begin{table}[H]
\centering
\caption{Especificaciones geométricas en superficies críticas de estanqueidad}
\label{tab:tolerancias_superficies}
\vspace{2mm}
\begin{tabularx}{\textwidth}{X r l}
\toprule
\rowcolor{categorybg}
\textbf{Característica Geométrica} & \textbf{Tolerancia Máx.} & \textbf{Estándar de Referencia} \\
\midrule
Planitud de Cara de Sellado & $0.05\,\text{mm}$ & ISO 1101 (Clase A) \\
\rowcolor{rowalt}
Paralelismo Tapa--Base & $0.03\,\text{mm}$ & ISO 1101 (Clase A) \\
Concentricidad de Torretas M2 & $0.02\,\text{mm}$ & ISO 1101 (Clase A) \\
\rowcolor{rowalt}
Rugosidad Superficial en Canal & $R_a \le 1.6\,\mu\text{m}$ & ISO 4287 \\
\bottomrule
\end{tabularx}
\end{table}

\section{Ubicación e Integración de Antenas}

\subsection{Especificaciones Electromecánicas}

El BOM define tres antenas independientes, una por banda; ninguna cubre más de un servicio:

\begin{itemize}[leftmargin=*]
    \item \textbf{LoRa (915 MHz)}: antena flexible de $8\,\text{cm}$ con conector u.FL (monopolo de cuarto de onda; el listado del proveedor indica $2\,\text{dBi}$ y no tiene MPN). Su rendimiento depende del plano de masa de la PCB y el alcance real debe probarse. La señal sale del pad ANT del módulo DX-LR30-900M22S por una línea de $50\,\Omega$ en la PCB hasta el conector u.FL.
    \item \textbf{Wi-Fi/BLE (2.4 GHz)}: antena FPC Würth 2015301200 conectada por u.FL al módulo ESP32-S3-WROOM-1U-N16R8. Su ganancia no está verificada: no hay hoja de datos en el repositorio.
    \item \textbf{GNSS (1575.42 MHz)}: parche cerámico pasivo de $18\times18\times4\,\text{mm}$, con MPN pendiente. El MAX-M10S integra LNA y filtro SAW; el plano de masa del parche se define en el diseño de la PCB.
\end{itemize}

\subsection{Criterios de Posicionamiento Radiativo}

La posición de cada antena en la carcasa y en la PCB está \textbf{POR DECIDIR}. Los siguientes márgenes de aislamiento electromagnético son criterios de diseño propuestos, aún sin validar, para aplicar a cada antena una vez decididas sus posiciones, junto con un aislamiento suficiente entre las bandas de $915\,\text{MHz}$, $1575\,\text{MHz}$ y $2.4\,\text{GHz}$:

\begin{itemize}[leftmargin=*]
    \item Separación superior a $10.0\,\text{mm}$ respecto a los bordes laterales del PCBA.
    \item Orientación paralela al eje longitudinal para maximizar la polarización vertical.
    \item Distancia mínima de $5.0\,\text{mm}$ a elementos metálicos (batería LiPo y masa del panel solar).
\end{itemize}

\subsection{Fijación Mecánica}

Aplica a la antena FPC de $2.4\,\text{GHz}$:

\begin{itemize}[leftmargin=*]
    \item \textbf{Adhesivo Industrial}: Lámina acrílica transferible $3\text{M}\,467\text{MP}$ de alta adherencia.
    \item \textbf{Protección Mecánica}: Film superior de poliimida (Kapton) de $50\,\mu\text{m}$.
    \item \textbf{Radio de Curvatura Mínimo}: $R_{\text{mín}} = 2.0\,\text{mm}$ en el cuello de conexión U.FL.
\end{itemize}

\section{Integración del Panel Solar y Carga MPPT}

\subsection{Parámetros Electromecánicos del Panel}

\begin{itemize}[leftmargin=*]
    \item \textbf{Tecnología}: Células fotovoltaicas monocristalinas de alta eficiencia ($18\text{--}20\%$).
    \item \textbf{Tensión Nominal}: $6\,\text{V}$ y potencia de al menos $1\,\text{W}$. El MPN del panel está \textbf{POR DECIDIR}: el HTJ-2W-D del BOM no declara su tensión.
    \item \textbf{Tensión en Vacío ($V_{oc}$)}: por definir según el panel elegido, claramente mayor que $6\,\text{V}$. El CN3791 con su divisor MPPT para $6\,\text{V}$ solo inicia la carga por encima de unos $6.17\,\text{V}$ en VCC, de modo que un panel de $5.5\,\text{V}$ en vacío nunca la iniciaría.
    \item \textbf{Corriente de Cortocircuito ($I_{sc}$)}: por definir según el panel elegido.
\end{itemize}

\subsection{Encapsulado y Orientación}

El panel está integrado en la cara superior expuesta del collar. La superficie activa cuenta con un encapsulado de fluoropolímero ETFE con tratamiento anti-UV diseñado para un objetivo de más de $5000\,\text{horas}$ de exposición solar continua sin amarillamiento (pendiente de ensayo).

\section{Materiales, Fijaciones y Acabados}

\subsection{Carcasa Externa}

\begin{itemize}[leftmargin=*]
    \item \textbf{Material Base}: Policarbonato con $20\%$ de fibra de vidrio (\text{PC+GF20\%}), resina grado UV.
    \item \textbf{Color de Carcasa}: Negro mate estabilizado contra la radiación solar.
    \item \textbf{Texturizado}: Acabado SPI-C1 / VDI 24 ($R_a = 3.2\text{--}6.3\,\mu\text{m}$) antideslizante.
\end{itemize}

\subsection{Tornillería y Torques de Ensamble}

\begin{itemize}[leftmargin=*]
    \item \textbf{Tornillos de Cierre}: Tornillos M2x6 de cabeza allen con inserto roscado M2 (referencias del BOM: PEM 50017102 y PEM-SL2-02).
    \item \textbf{Material de Fijación}: Acero inoxidable A2 (AISI 304) pasivado.
    \item \textbf{Torque de Apriete Controlado}: $0.6\text{--}0.8\,\text{N}\cdot\text{m}$ para evitar fisuras en torretas.
\end{itemize}

\section{Conclusiones}

El diseño mecánico del collar DR-VF-01 satisface los requerimientos de operabilidad en campo abierto ganadero. El uso de la resina PC+GF20\% combinado con la junta tórica EPDM en compresión controlada asegura la estanqueidad IP67. Asimismo, el diseño prevé tres antenas independientes (LoRa, Wi-Fi/BLE y GNSS) y el panel solar; sus posiciones en la carcasa están \textbf{POR DECIDIR}, y la transmisión RF sin apantallamiento y la recarga solar continua quedan por verificar.